%%%%%%%%%%%%%%%%%%%%%%%%%%%%%%%%%%%%%%%%%%%%%%%%%%%%%%%%%%%%%%%%
% Master's Thesis -- Colorado State University
% Migrated to the CSU thesis template (thesis.cls by Elliott Forney)
%%%%%%%%%%%%%%%%%%%%%%%%%%%%%%%%%%%%%%%%%%%%%%%%%%%%%%%%%%%%%%%%

\documentclass[master]{thesis}

% ams math packages
\usepackage[cmex10]{amsmath}
\usepackage{amsthm,amssymb}

% graphics
\usepackage[pdftex]{graphicx}
\graphicspath{{./figures/}}
\DeclareGraphicsExtensions{.pdf,.png,.jpg}

% subfigures
\usepackage[caption=false]{subfig}

% book-quality tables
\usepackage{booktabs}

% urls and links
\usepackage{url}
\urlstyle{same}

% math/engineering style citations
\usepackage{cite}

% custom reference commands provided by the template
\newcommand{\eref}[1]{\eqref{#1}}        % cite equation
\newcommand{\fref}[1]{Figure~\ref{#1}}   % cite figure
\newcommand{\cref}[1]{Chapter~\ref{#1}}  % cite chapter
\newcommand{\sref}[1]{Section~\ref{#1}}  % cite section/sub(sub)section
\newcommand{\aref}[1]{Appendix~\ref{#1}} % cite appendix
\newcommand{\tref}[1]{Table~\ref{#1}}    % cite table

% Title Page metadata
%%%%%%%%%%%%%%%%%%%%%%%%%%%%%%%%%%%%%%%%%%%%%%%%%%%%%%%%%%%%%%%%

\title{Answer-Structure Prompting for Small Language Models}

\author{Eric Redman}
\email{C830573395@colostate.edu}
\department{Department of Computer Science}
\semester{Fall 2026}

\advisor{Nikhil Krishnaswamy}
\committee{Anura P. Jayasumana}
\committee{Ravi Mangal}

% Copyright Page
%%%%%%%%%%%%%%%%%%%%%%%%%%%%%%%%%%%%%%%%%%%%%%%%%%%%%%%%%%%%%%%%

\mycopyright{%
Copyright by Eric Redman 2026 \\
All Rights Reserved
}

% Abstract
%%%%%%%%%%%%%%%%%%%%%%%%%%%%%%%%%%%%%%%%%%%%%%%%%%%%%%%%%%%%%%%%

\abstract{%
Small language models are limited by capacity, and prompting is one way to compensate at inference time. This work investigates whether supplying a small model with an answer structure matched to the question, a schema naming the components a valid answer must contain, improves its accuracy on extractive question answering. A taxonomy of answer structures was derived from SQuAD and examined through two studies. A six-model study spanned 1B to 8B parameters. A fine-tuning study varied how the structure was trained and delivered. The matched structure did not improve accuracy and in most cases reduced it, while the source passage improved every model. The failure lay in the form of delivery rather than in the structure itself. The same structure expressed as prose improved accuracy where a schema did not. While the result is bounded to extractive question answering, the broader takeaway is that structure can help a small model, provided it arrives in a form the model was trained to read.
}

% Acknowledgements
%%%%%%%%%%%%%%%%%%%%%%%%%%%%%%%%%%%%%%%%%%%%%%%%%%%%%%%%%%%%%%%%

\acknowledgements{%
Thank you to my advisor, Nikhil Krishnaswamy, for his guidance and help.
}

% Use of Generative AI disclosure
%%%%%%%%%%%%%%%%%%%%%%%%%%%%%%%%%%%%%%%%%%%%%%%%%%%%%%%%%%%%%%%%

\newcommand{\aidisclosure}{%
This thesis made use of generative AI tools, disclosed here in accordance with Colorado State University policy.

\vspace{1em}

Claude Haiku 4.5 was used as an automated annotator to label questions during the derivation of the answer-structure taxonomy, as described in the Methodology.

\vspace{1em}

Kiro, an AI coding assistant, implemented the experiment code from the author's design and assisted in revising and drafting portions of the text. It was also used to help verify citations and to check reported results against the underlying data. All edits, changes, and additions were directed and reviewed by the author. The research design, analysis, results, and conclusions are the author's own, and no data, experimental results, or citations were fabricated.
}

\newcommand{\makeaidisclosure}{%
    \prelimtocentry{Use of Generative AI}

    \begin{flatcenter}
        USE OF GENERATIVE AI

        \vspace{3em}
    \end{flatcenter}

    \aidisclosure \par

    \newpage
}

\begin{document}

\frontmatter

\maketitle
\makemycopyright
\makeabstract
\makeacknowledgements
\makeaidisclosure

\tableofcontents
\listoftables
\listoffigures

\mainmatter

% ============================================================
% BODY -- ported from main.tex, section levels shifted:
%   \section    -> \chapter
%   \subsection -> \section
% ============================================================

\chapter{Introduction}
\label{ch:intro}
% ============================================================

The word \emph{structure} recurs throughout this work in three distinct senses, and separating them at the outset prevents the ambiguity that would otherwise follow. The first is the explicit schema supplied to the model in its input, a template naming the components a valid answer must contain, provided before generation begins; this is the \textbf{answer structure}, the object this work supplies and tests. The second is the organization a model holds internally, the patterns for arranging an answer that are encoded in its parameters during training and retained at a resolution set by scale; this is \textbf{structural encoding}. The third is the relational and syntactic organization already present in well-formed language itself, independent of any schema; this is the \textbf{structure of prose}. Each later use of the term refers to one of these three.

Large language models have proven capable across a broad range of tasks. This capability comes at a cost in scale and resources, which concentrates their development in the hands of a few \cite{eloundou2023gpts, strubell2019energy}. Smaller models are far more accessible, but they are weaker, particularly on general tasks \cite{smalllmsready2024, slm_survey2024}. The usual explanation is scale: with fewer parameters, a model retains less of what it saw in training \cite{kaplan2020scaling, costofdownscaling2023}. Their language abilities still survive the reduction \cite{linguistic_competence_scale2025, emergent_reducedscale2024}. What a small model lacks is therefore something other than language. Whether it can be supplied without adding parameters is the central question of this work.

A small model may lack not language but structure. Such a structure can be supplied at inference time from outside the model, rather than built in through training. This is the possibility the present work investigates. The hypothesis is that an answer structure matched to the question, a schema naming the components a valid answer must contain, narrows the model's prediction space and produces a more accurate answer. The approach would in principle apply to any model and any task whose answer structure is known.

To test this, a taxonomy of answer structures was derived from the Stanford Question Answering Dataset and examined through two studies. A six-model study, spanning 1B to 8B parameters, supplied each model with the matched structure, a random structure, and the source passage across conditions. A fine-tuning study, on a single model, varied how the structure was trained and how it was delivered.

The matched structure did not help. It failed to improve any model and in most cases lowered accuracy, while the source passage improved every one. The failure lay in the form of delivery rather than in the structure itself. The same structure expressed as prose improved accuracy where the schema did not. A small model does not reject structure. It reads structure only in the form it was trained on, and a rigid schema is not that form. This result comes from extractive question answering, a setting least favorable to the hypothesis. Here answers are short spans with almost no form for a structure to organize. That the effect is negative even here is informative, and whether it holds where answers must be composed is a question left open.

This work makes three contributions. The first concerns delivery. It identifies the form in which a structure is supplied as the factor governing whether that structure helps a small model. An imposed schema fails where a prose rendering of the same structure succeeds. The second is the negative result itself. Across six models and a fine-tuning study, an imposed answer structure did not improve accuracy, and the fine-tuning study reports a confidence interval for every comparison. The third is an explanation. A small model represents language as a learned distribution over ordinary text, and it can use a structure only when that structure arrives in the same form.

% ============================================================
\chapter{Related Work}
\label{ch:lit}
% ============================================================

\section{Large Language Models}
\label{sec:llm}

Large language models are a class of deep neural networks built on the transformer architecture and trained on massive corpora of text \cite{vaswani2017attention, llmoverview2023}. The transformer encodes tokens, discrete units representing words or subwords. Through its self-attention mechanism, every token attends to every other token in the sequence \cite{vaswani2017attention}, allowing the model to encode the structural features of language. These include the relationships between words and phrases, syntax, and semantics. The same mechanism allows the model to condition on provided context, including prompts and scaffolds supplied at inference time \cite{brown2020gpt3}. Parametric capacity is the key constraint on what a model can encode from the training corpus \cite{kaplan2020scaling}, but regardless of size, the fundamental mechanism remains \cite{costofdownscaling2023}.

By design, a large language model is a prediction engine. It does not store exact or factual information, but the statistical relationships between tokens learned during training. When a model receives a prompt, it performs a forward pass that produces a probability distribution over the full vocabulary. Every possible continuation receives a likelihood. The output has no guarantee of correctness. It is simply the most statistically probable continuation given the input and the model's learned parameters. Furthermore, the model produces an identical distribution for a given input every time. Any difference in output depends on the sampling configuration. The result is a prediction shaped entirely by the input it receives.

Once generation begins, each selected token becomes the input for the next prediction, narrowing the output distribution further \cite{brown2020gpt3}. A wrong prediction at any point can alter the trajectory of the entire response \cite{exposure_bias_ilp2022, lpsr2026}. This is the nature of autoregressive decoding: the model generates one token at a time, each conditioned only on what came before, with no mechanism for lookahead or self-correction \cite{brown2020gpt3, exposure_bias_ilp2022}. The quality of the output is therefore contingent on the initial prompt \cite{prompt_control_theory2023}. This may not hold for all large language models, but it is the confirmed architecture of every open model evaluated in this paper.\footnote{Closed-source models may employ internal techniques that partially mitigate autoregressive limitations, but their architectures cannot be independently verified \cite{liesenfeld2024closed}.}

What the model receives as input directly determines the distribution from which tokens are selected.

\section{Scale and Structure}
\label{sec:scale}

Scale, as used in the literature, refers to the number of parameters in a language model \cite{kaplan2020scaling}. It is the dominant predictor of general performance. Larger models consistently outperform their smaller counterparts across a wide range of domains \cite{kaplan2020scaling, hoffmann2022chinchilla}. The prevailing explanation attributes this to capacity: a larger parameter count encodes more of the information in the training data, and smaller models are forced to compress or discard a portion of it \cite{llmoverview2023, knowledgecapacity2024}.

This can be observed in Brown et al. \cite{brown2020gpt3}. A GPT-3 model with 175B parameters executed few-shot tasks that GPT-3 at 1.3B parameters was incapable of performing. There were no changes to the architecture or training regime. The only variable was the number of parameters.

Scale alone, however, does not fully explain the difference. Even the largest models fail on multi-step reasoning tasks and produce hallucinations \cite{math_reasoning_limits2024, lrm_complexity2025}. The gap is conditional, not absolute. On specific tasks in specific domains, small models given appropriate context can match and even exceed large ones \cite{smalllmsready2024, modelsizematters2024}.

Architecturally, small models are not distinct from larger ones. They share the same generation mechanisms and are trained on the same corpora toward the same target distribution \cite{godey2024softmax, costofdownscaling2023}. The difference in scale, then, must be the resolution of the approximation stored in the parameters.

The advantage provided by scale is structural. The capacity of larger models allows them to encode the same structures at higher resolution \cite{godey2024softmax}. Consider an image file: reduce the pixels and the image remains, but in a degraded form. Reduce it too far and the image breaks down entirely. In this analogy, the image represents the structural patterns of language. These include how to format a response, how to relate concepts across a sequence, and how to organize an answer to match the type of question being asked. Yet unlike a deterministic structure such as a JSON schema, in an LLM these patterns are statistically encoded and non-deterministic, resulting in possible corruption regardless of model size \cite{math_reasoning_limits2024, lrm_complexity2025}.

Prompting techniques are effective because they compensate for degraded structural encoding. Supplying a small model with a structured, task-matched guide at inference time improves its performance without any change to its weights \cite{han2025blueprints}. This demonstrates that structure, in addition to capacity, can have an impact on output quality.

Prompting alone does not solve generalization. Whether a small model can compensate for the parametric constraint at inference time remains open.

\section{Fine-Tuning and Parameter-Efficient Adaptation}
\label{sec:finetune-lit}

Often the first step to improving a model is affecting its weights \cite{peft_survey2024}. The logic is simple: the parameters store encoded information, and by updating this encoding to better match a given task, the model should improve.

Full fine-tuning updates every parameter; adapter methods insert small trainable modules between existing layers; prefix tuning prepends learnable tokens to the input \cite{hu2022lora, houlsby2019adapters, li2021prefixtuning}. Some, like LoRA, update smaller subsets of parameters to reduce computational cost \cite{hu2022lora}. Each updates a model's learned distribution by reshaping its weights on task-specific data.

Instruction tuning deserves particular attention. This method of fine-tuning relies on instruction-following data to teach the model how to respond to a variety of task prompts, and has been shown to improve generalization to tasks unseen during tuning \cite{instructiontuning2022}. Fine-tuning applied to a broad aspect of model behavior, rather than to a single narrow end task, can raise performance across many downstream tasks at once. The effect has limits. As instruction-tuning data grows, performance on knowledge and reasoning benchmarks declines \cite{alignment_tax2024}. The gain in breadth is real, but it is purchased against a capacity that remains fixed.

A related approach worth highlighting is knowledge distillation. This process trains a smaller model using a larger model's outputs, updating the student's weights with the teacher's output encodings \cite{hinton2015distilling}. As the teacher has a higher capacity for encoding, its outputs are of higher quality, resulting in more focused training than general fine-tuning. The student model, however, is still bounded by its own capacity. No matter how focused the training outputs are, the encoding is constrained by the same dimensional bottleneck (\sref{sec:scale}) \cite{godey2024softmax, costofdownscaling2023}. Distillation shifts where a fixed capacity is spent rather than expanding it.

Modifying a model's parameters carries a further cost. Fine-tuning narrows a model's probability distribution toward the training task, which can degrade performance on others \cite{catastrophic_forgetting2023}. When applied to a broad aspect such as instruction-following, this narrowing is the mechanism of the gain; when applied to a single narrow task, the same mechanism can erode breadth. Fine-tuning is therefore a tradeoff to be managed rather than a lever that only ever helps or only ever harms. It reshapes how a fixed capacity is allocated. A complementary path leaves the weights untouched and instead supplies what the parameters do not reliably encode at inference time.


\section{Prompt Engineering and In-Context Learning}
\label{sec:prompt}

Prompt engineering shapes model output by constructing the input at inference time. Additional context aligned with the input query can narrow the model's prediction space, reducing output error \cite{whendoescontexthelp2026}. These techniques do not update model parameters and are agnostic to model version.

Several variations of prompting techniques exist. Zero-shot, for example, is the most basic, providing only the initial question and a framing of the task \cite{brown2020gpt3}. This relies heavily on the model's parametric knowledge with minimal contextual constraint. Few-shot techniques provide more context in the form of example input-output pairs. Providing these concrete patterns to the model helps constrain both the format and content of its response \cite{brown2020gpt3, sahoo2024promptsurvey}. Techniques like chain-of-thought provide explicit context related to the processing of the prompt, directing the model to reason through intermediate steps before producing an answer \cite{wei2022cot}. The common thread is that all these techniques provide additional context to the input. In addition to activating more of the model's encoded relationships, the explicit structure between tokens in the prompt becomes additional signal the model can attend over \cite{olsson2022inductionheads}, narrowing the prediction space. This mechanism is known as in-context learning (ICL) \cite{brown2020gpt3}.

ICL operates through a model's attention rather than its parametric capacity, and the structure it provides is external to the model, allowing it to function when capacity is limited. This makes ICL scale-resistant \cite{costofdownscaling2023}. As an input-side intervention, ICL can also be applied across a variety of domains and models without prior training or modification \cite{brown2020gpt3}. Few-shot prompting, for instance, does not prescribe what the examples must be, only that examples improve performance \cite{sahoo2024promptsurvey}. The prompt structure remains constant while the content varies with the task.

This does not mean prompt engineering alone can solve the problem of generalization at a small scale. The existing techniques work well, but still require the model to produce outputs based on information encoded in its parameters. When the model has fewer parameters, that information suffers the same resolution loss established in \sref{sec:scale}.

Chain-of-thought prompting highlights this explicitly. It works well on larger models \cite{wei2022cot} but suffers significant degradation on smaller ones \cite{cot_degradation_slm2025, smallmodel_learnability2025}. Even with the additional structure provided by the prompt, iteratively generating on an incorrect encoding causes the prediction space to expand \cite{exposure_bias_ilp2022, lpsr2026}.

Prompt engineering effectively narrows the prediction space, but still cannot entirely account for the parametric capacity gap. A consistent pattern across these techniques is that more specific, better-structured input yields a sharper distribution.


\section{Retrieval-Augmented Generation}
\label{sec:rag}


Retrieval-Augmented Generation, as the name implies, augments generation by retrieving additional context to construct the prompt at inference time \cite{lewis2020rag, rag_survey2023}. This retrieval pulls information from an external source.

Like prompt engineering (\sref{sec:prompt}), RAG uses in-context learning to improve a model's performance without modifying its internal weights \cite{lewis2020rag}. The difference is in what each provides. Prompt engineering supplies structure and framing designed by the author to activate relationships the model might not otherwise prioritize. RAG supplies specific content retrieved per-query to cover information the model's parameters may not encode. The two are not mutually exclusive; RAG can augment prompt engineering and vice versa \cite{prompt_eng_rag_slm2026}.

Retrieval-Augmented Generation's value is in specificity. Retrieval from an external source supplies explicitly task-relevant content at inference time \cite{lewis2020rag}. These can range from factual content pulled from a static knowledge base to passages selected dynamically based on query similarity \cite{rag_survey2023}. As established in \sref{sec:prompt}, providing more relevant context in the prompt improves performance \cite{brown2020gpt3}. This additional context increases the implicit relational structure during generation, reducing the prediction space. Certain RAG techniques provide organization beyond just relevant context. For example, question-decomposition approaches break a complex query into sub-questions and retrieve separately for each, providing more precise content at each step and a clear breakdown of a given problem \cite{question_decomp_rag2025}. As with prompting (\sref{sec:prompt}), more context improves performance, more context with structure improves performance further, and context with both structure and specificity outperforms either alone.

As an inference-time technique, RAG is both generalizable and scale-resistant. A retrieval system can be applied to any model, trained or untrained, and depending on the technique, applied across a variety of domains without modification \cite{rag_survey2023}. Its core mechanism operates through attention rather than parametric capacity, meaning benefits persist when models are smaller \cite{costofdownscaling2023}. However, the same limitation identified in prompt engineering (\sref{sec:prompt}) applies: there remains a reliance on the model's encoded information to organize its output. RAG supplies the content an answer needs, but the burden of shaping that content into a well-formed answer still falls on the model.

\section{Unconstrained Prediction}
\label{sec:unconstrained}


When a model does not produce the expected or correct answer for a given task, it is common to call the error a hallucination \cite{ji2023hallucination, hallucination_unified2024}. At the architectural level, this is a misleading description. The model did not make a mistake nor imagine a false reality. It simply produced an output that did not match the task. As introduced in \sref{sec:llm}, for every forward pass the model produces a probability distribution derived entirely from the prompt \cite{brown2020gpt3, llmreasoning_stochastic2024}. These errors arise from the ordinary statistical behavior of fitting the training distribution \cite{whyllmshallucinate2024}. There is no alternative failure mode.

LLMs use temperature to reshape the softmax probability curve and a random seed to select from that distribution, creating the appearance of variety \cite{temperature_impact2025}. Like the neural networks they are built on, if these values are static, LLMs will always produce the same output. When these values vary, they produce different but adjacent outputs. Seed determines which token is drawn from the distribution; the shape of that distribution is unchanged, so divergence is bounded \cite{seed_temperature_stability2025}. Temperature reshapes the distribution itself. Selecting adjacent tokens in the probability space can cause the output to diverge more significantly \cite{temperature_impact2025}. Neither actively improves the generation.

As established in \sref{sec:llm}, the input defines the trajectory of the prediction, making it a primary point of intervention. If the quality of the input determines the quality of the output, then it is critical the input has the right context. Chain-of-thought demonstrates what happens when this context is absent. When the input does not sufficiently constrain the trajectory, the model generates purely from its internal representations. If those representations do not match the task space, the output diverges significantly \cite{cot_degradation_slm2025, smallmodel_learnability2025}.

Larger models often have this specificity and structure encoded in their parameters. Smaller models lack this resolution (\sref{sec:scale}). Prompt engineering can supply structure (\sref{sec:prompt}), but the structure it provides is generic and is applied uniformly regardless of the task. RAG supplies specificity (\sref{sec:rag}), but not in a form adapted to the specific task. Even combined, these techniques leave the model to organize its output according to the specific task.

\section{Question Classification in QA}
\label{sec:qclass}


Question classification is well-established in computer science and foundational to question answering systems. Li and Roth demonstrate that a taxonomy of question classes exists and that these classes can be identified from the question text alone \cite{liroth2002questions}. Different classes of questions have different associated features. For example, some are factual, others map relations. Bolotova et al. establish a taxonomy of this kind for non-factoid questions, whose answers are open-ended rather than short spans \cite{nonfactoid_taxonomy2022}. Its class signals are observable at the surface level and even extend across subject domains. If answer structure is finite and identifiable there, it follows that answer structure would be identifiable for the extractive, span-based questions studied here (\cref{ch:method}), whose answers are less complex.

If questions are classifiable, and every question has a relevant answer, there must be some structure or set of features that can be used to identify groups of answers. If questions are classifiable from their text \cite{liroth2002questions}, the question type may also indicate the form its answer takes. A factual question suggests a single fact. Analogical reasoning involves a relational mapping between two subjects \cite{analogical_reasoning2022}. Multi-step problems tend toward a sequence of intermediate conclusions \cite{wei2022cot}. These relationships are not equally explicit, but they are observable and learnable from examples \cite{devlin2019bert, xu2019qc_bert}.

Work in lexical semantics offers independent precedent for this idea. Meaning can be decomposed into a finite, repeatable structure rather than treated as surface form. Pustejovsky's Generative Lexicon represents word meaning through a small set of relational roles \cite{pustejovsky1991generative}. A qualia structure describes an object through roles capturing what it is made of, what distinguishes it, what it is for, and how it comes about. The noun \emph{book} is characterized in part by a telic role of reading and an agentive role of writing. A \emph{dictionary} shares the same formal role but differs in purpose, its telic role being reference. The roles are fixed while their values vary with the word. This framework concerns lexical meaning rather than question answering, but it establishes that the semantics of an expression can be captured by a bounded set of relational roles.

A taxonomy of answer forms would align with the properties described in \sref{sec:unconstrained}. The relational patterns are content-independent and apply across answers regardless of task. The dimensions within each form can be provided to a model before generation begins. A form derived from the question type is consistent across domains, since it depends on the question's structure rather than its subject. No such taxonomy of answer structures appears in the existing literature.

% ============================================================
\section{Similar Approaches}
% ============================================================

Several recent works cover adjacent problems, and each contributes an element of the present approach.

SMART proposes that properties of the answer can be identified from the input question alone \cite{setty2023smart, smart2021}. This is done by predicting the ontological entity type an expected answer requires. The predicted type is used to narrow the retrieval search space so that only entities of that type are considered. This work demonstrates a question contains enough signal to identify features of an answer before producing one.

Dynamic Template Selection classifies incoming prompts by semantic complexity to select a pre-built response template \cite{yadavalli2025dts}. The templates determine the output's verbosity and structure. The objective is to reduce the cost of generation through more efficient token usage. For the present purposes, this work highlights that classification of an input can map directly to a templated taxonomy.

TemplateRL constructs a library of problem-solving templates \cite{wu2025templaterl}. Like the answer structures described here, these templates provide structured context to assist in generation. The difference is that TemplateRL's templates provide a set of reasoning action sequences, indexed by problem complexity and retrievable per-problem. This work establishes that a structured library can be maintained and retrieved on a per-prompt basis to improve generation.

The Blueprints devised by Han et al. structure step-by-step reasoning to guide generation \cite{han2025blueprints}. By providing these as input tokens to a small model, they are shown to increase performance. No weight updates are required; the intervention is purely at inference time. This demonstrates that structured guidance provided as a prompt can improve small model performance without fine-tuning.

Constrained decoding masks illegal tokens at each step of the generation process \cite{willard2023guided}. This guarantees the output conforms to a specified grammar. It operates at decode time and limits the output distribution directly. Like an answer structure, constrained decoding illustrates that the distribution of a prediction can be sharpened with a structural constraint.

Each of these works establishes a piece of the approach studied here: that a question carries signal about its answer, that classification can map to a template, that a structured library can be retrieved per problem, and that inference-time structure can aid a small model. None combines them into task-matched answer structures evaluated for their effect on small-model accuracy.

% ============================================================
\chapter{Methodology}
\label{ch:method}
% ============================================================

\section{Research Questions}

The hypothesis is that a matched answer structure improves a small model's accuracy on extractive question answering. It is examined through seven research questions:

\begin{enumerate}
  \item \textbf{RQ1.} Does a matched structure improve accuracy over no structure?
  \item \textbf{RQ2.} Does the specific structure matter, or does any structure help equally?
  \item \textbf{RQ3.} Does the structure contribute independently of the source passage?
  \item \textbf{RQ4.} Is the result consistent across model sizes?
  \item \textbf{RQ5.} Does fine-tuning on the structure improve model effectiveness?
  \item \textbf{RQ6.} Does the form in which the structure is delivered change the result?
  \item \textbf{RQ7.} Can a structure carrying the full content rival prose carrying the same content?
\end{enumerate}

Two studies address these questions. The six-model study covers RQ1 through RQ4, and the fine-tuning study covers RQ5 through RQ7.

\section{The Answer Structure Hypothesis}
\label{sec:taxhyp}

An answer structure is a schema that specifies what a valid answer must contain without specifying what the answer is. Each entry in the schema is a dimension: a required component of the answer. The schema defines what kinds of information are needed, how they relate, and how they must be organized. It provides explicit relational context in the input, constraining the prediction space before generation begins. When prompted, the model's remaining task is to populate the structure with relevant content.

The following examples illustrate how question types produce distinct answer structures. A relational question such as ``How does A relate to B?'' requires specific dimensions: (1) the nature of A, (2) the nature of B, and (3) how they relate. A causal question such as ``Why did X happen?'' calls for a different set: (1) the event, (2) the preceding causes, and (3) the causal mechanism. A procedural question such as ``How do you do X?'' differs again: (1) the task, (2) the steps required, and (3) the order of the steps. In each case the answer requires subjects, a mechanism of relation, and joining details. The dimensions vary across types, but the structural pattern is consistent, and the number of dimensions is not fixed. In all cases, the structure can be identified from the question and provided before generation begins. These examples are illustrative, not exhaustive.

\begin{figure}[ht]
\centering
\begin{minipage}[t]{0.45\textwidth}
\textbf{Answer Structure (provided):}
\begin{verbatim}
{
  "event": "",
  "causes": [],
  "mechanism": ""
}
\end{verbatim}
\end{minipage}
\hfill
\begin{minipage}[t]{0.45\textwidth}
\textbf{Populated by model:}
\begin{verbatim}
{
  "event": "sinking of the Titanic",
  "causes": [
    "ship enters iceberg territory",
    "collision occurs",
    "hull is punctured"
  ],
  "mechanism": "causal chain"
}
\end{verbatim}
\end{minipage}
\caption{An answer structure for the question ``Why did the Titanic sink?'' The left shows the empty structure provided to the model before generation. The right shows the structure populated by the model.}
\label{fig:answer_structure_example}
\end{figure}

An answer structure supplies additional context that narrows the prediction space through in-context learning (\sref{sec:prompt}). The structure contributes dimensional hooks, such as subjects or mechanisms, as input tokens; the more of these are present, the more attention is constrained and the sharper the resulting distribution. For a smaller model, that sharpened distribution compensates for limited parametric resolution (\sref{sec:scale}). The mechanics of how a scaffold delivers this at inference time are taken up in Prompt Scaffolding below.

Answer structures also support generalization. As a defined taxonomy, they apply across many problems with the same process: identify the question type, provide the structure, populate the dimensions. The only change between tasks is the structure itself. Unlike few-shot prompting, which provides examples a model must infer a pattern from, an answer structure provides a direct mapping to the answer form. Theoretical work suggests a broader principle: a prompt matched to a task's structure outperforms a generic one \cite{taskspecific_prompts2025}. The taxonomy developed here applies that principle at the level of individual question prompts.

\section{Dataset}

The dataset used in this paper is the Stanford Question Answering Dataset (SQuAD) \cite{rajpurkar2016squad}. SQuAD consists of extractive question answering over Wikipedia passages and is a widely used benchmark for automated QA evaluation. In an extractive task the answer is a short span drawn directly from the passage, so answers carry little internal structure of their own. This isolates the effect of the imposed structure, testing it where the answer contributes least. Any findings are scoped to this setting rather than to tasks with longer, compositional answers. Question answering as a format spans many structural types such as factual, causal, temporal, and descriptive. One format covering many distinct answer forms makes it suitable for empirical derivation of a taxonomy. Gold-standard span answers further enable automated scoring via Exact Match and F1 without requiring human judgment. SQuAD also provides a source passage from which the answer is extracted. The passage serves as the context variable in the experimental design: it can be included or excluded, allowing the study to test whether the answer structure contributes independently of source context (RQ3). Extension to other domains requires only re-derivation of the taxonomy.

\section{Taxonomy Derivation}

The taxonomy defines the set of prompt scaffolds. It contains 10 classes of answer structure, each representing a distinct relational role the answer plays relative to the question. Each class maps to a unique answer structure consisting of a set of labeled dimensions. Classes are defined by the structural form of the answer, not by surface features of the question such as the interrogative. The set is static and predefined prior to inference.

The 10 classes are:

\begin{enumerate}
  \item \textbf{Object Identification} --- the answer names or identifies a concrete object, entity, document, technology, or substance.
  \item \textbf{Concept/Person Identification} --- the answer identifies an abstract concept, classification, or person by a defining property.
  \item \textbf{Multipart Identification} --- the answer requires identifying multiple distinct components.
  \item \textbf{Quantitative} --- the answer is a number: a count, measurement, proportion, or recorded statistic.
  \item \textbf{Temporal} --- the answer is a time expression: a date, year, era, duration, or range.
  \item \textbf{Place} --- the answer is a geographic location or spatial reference.
  \item \textbf{Agent} --- the answer is who or what performed an action.
  \item \textbf{Descriptive} --- the answer describes a quality, property, or characteristic of something.
  \item \textbf{Causal} --- the answer explains why something happened: the cause, reason, or enabling mechanism.
  \item \textbf{Resultative} --- the answer is the outcome, effect, or product of an action or event.
\end{enumerate}

Each class maps to a schema that specifies the dimensions a valid answer must address. Every schema contains two common pre-filled fields: \texttt{answer\_type}, which indicates cardinality (single or multiple), and \texttt{answer}, which names the kind of answer expected. The remaining fields are class-specific and empty, defining the structural components the model must populate. For example, the Causal class maps to a schema containing the effect, the cause, and the mechanism linking them. The Agent class maps to a schema containing the agent, the action performed, and the object of that action. The schemas differ in both the number and nature of their dimensions, providing structurally distinct guidance per class.

\begin{figure}[ht]
\centering
\begin{minipage}[t]{0.30\textwidth}
\textbf{Temporal:}
\begin{verbatim}
{
  "answer_type": "single",
  "answer": "time",
  "event": "",
  "time_expression": "",
  "temporal_type": ""
}
\end{verbatim}
\end{minipage}
\hfill
\begin{minipage}[t]{0.30\textwidth}
\textbf{Causal:}
\begin{verbatim}
{
  "answer_type": "single",
  "answer": "cause",
  "effect": "",
  "cause": "",
  "mechanism": ""
}
\end{verbatim}
\end{minipage}
\hfill
\begin{minipage}[t]{0.30\textwidth}
\textbf{Resultative:}
\begin{verbatim}
{
  "answer_type": "single",
  "answer": "outcome",
  "action": "",
  "actor": "",
  "outcome": ""
}
\end{verbatim}
\end{minipage}
\caption{Answer structure schemas for three taxonomy classes. Each schema contains two pre-filled fields (\texttt{answer\_type} and \texttt{answer}) that immediately signal what kind of answer is expected, followed by empty dimensions that define the structural components the answer must address.}
\label{fig:taxonomy_schemas}
\end{figure}

The taxonomy was derived empirically from the SQuAD dataset rather than from existing literature. Each question was described by a set of features, and the feature combinations that recurred across the dataset defined the classes. An initial four features were used: the interrogative, cardinality, the entity type, and the grammatical form of the answer. These proved insufficient. Questions with identical surface features often required different answer forms. A fifth feature, the relational context between the answer and the question, was added to capture the role the answer plays. For example, ``Who founded X?'' and ``Who was impacted by X?'' share the same interrogative and entity type but require different answer structures: the first identifies an agent, the second a resultative. This relational context proved to be the primary differentiator between classes.

The feature labeling was performed by a separate large language model acting as an annotator, distinct from the generator models later evaluated. The annotator was Claude Haiku 4.5, selected as a low-cost model adequate for structured labeling. It was given a rule set, established by inspection of the SQuAD development set, defining each of the five features and its permitted values. These were the six interrogatives; singular or multipart cardinality; six entity types; four grammatical answer forms; and eleven relational-context roles. The annotator was accessed through the Anthropic API using the model snapshot \texttt{claude-haiku-4-5-20251001}, with the rule set supplied as a fixed system prompt and no sampling parameters altered from their defaults. Questions were sent in batches of twenty. For each question the annotator received the question text and its gold answer and returned a value for each feature as structured output. The gold answer informed only this offline derivation and is not used at inference time. All 5,928 answerable questions in the SQuAD development set were labeled this way, each independently. The consolidation step served as implicit validation. Incoherent labels would not produce structurally meaningful clusters. The resulting output produced 254 unique feature combinations, 67 of which contained ten or more questions.

The 67 classes were reviewed and reduced by two criteria: the structural distinctiveness of the resulting answer schema, and class balance. Classes that mapped to identical schemas were merged, and classes too small to be represented reliably were absorbed into their nearest structural neighbor. Balancing the classes produces more complete and learnable training data. This produced 10 final classes, each with a unique schema and a largest-to-smallest class ratio of approximately 3:1. For example, temporal point and temporal duration both require the same dimensions (event, time expression, temporal type), differing only in content, not in guidance to the model. This taxonomy is scoped to SQuAD and is not claimed to be universal. It serves the purpose of demonstrating that a taxonomy of answer structures can be derived and operationalized.

Although derived independently, several of the resulting classes align with roles from an established descriptive framework. As discussed in \sref{sec:qclass}, the Generative Lexicon organizes meaning around a small set of relational roles \cite{pustejovsky1991generative}. The causal and agent classes echo its agentive role, which the Generative Lexicon illustrates with a causal chain. The multipart class echoes the constitutive role of components. Additional parallels can be found between the descriptive class and the formal role of distinguishing properties, and between the per-class dimensions and an argument structure.

\section{Prompt Scaffolding}

A scaffold is the input constructed for the model at inference time. Its purpose is to provide a structure that joins associated dimensional tokens into a single prompt. These tokens constrain attention and narrow the prediction space before generation begins (\sref{sec:taxhyp}). Each scaffold consists of up to three components: the answer structure schema, the input question, and an optional passage providing source context. The configuration is variable: depending on the condition, the answer structure and passage can each be included or excluded, while the question is always provided. The specific configurations are detailed in the study sections that follow.

The scaffold operates through in-context learning. The structure it supplies is external to the model's parameters, activating relationships through attention rather than adding to what the model has stored (\sref{sec:prompt}). In the six-model study the models receive no weight modification, so any effect is attributable to the scaffold alone. The fine-tuning study additionally adapts the model's weights, testing whether training changes how the scaffold is used.

\section{The Six-Model Study}
\label{sec:sixmodel-method}

The six-model study addresses RQ1 through RQ4. It evaluates six models across six conditions, each supplying the question with some combination of an answer structure and a source passage:

\begin{enumerate}
  \item \textbf{C1.} Question only
  \item \textbf{C2.} Question + passage
  \item \textbf{C3.} Question + matched structure
  \item \textbf{C4.} Question + passage + matched structure
  \item \textbf{C5.} Question + random structure
  \item \textbf{C6.} Question + passage + random structure
\end{enumerate}

The matched structure is the schema for the question's own class. The random structure is drawn from one of the nine classes that do not match it. C1 is the baseline for this study. It measures each model on zero-shot prompts containing only the question. Every other condition is compared against this baseline. RQ1 compares the matched conditions against their baselines, C3 against C1 without a passage and C4 against C2 with a passage. RQ2 compares the matched conditions against the random ones, C3 against C5 and C4 against C6. RQ3 asks whether structure adds anything to the passage, comparing the combined C4 against the passage-only C2. RQ4 compares these differences across the six models. \fref{fig:six_conditions} illustrates the six conditions for a single question.

\begin{figure}[ht]
\centering

\begin{minipage}[t]{0.30\textwidth}
\textbf{C1 (question only):}
\begin{verbatim}
In what country is
Normandy...?
\end{verbatim}
\end{minipage}
\hfill
\begin{minipage}[t]{0.30\textwidth}
\textbf{C2 (passage):}
\begin{verbatim}
The Normans gave
their name...
In what country is
Normandy...?
\end{verbatim}
\end{minipage}
\hfill
\begin{minipage}[t]{0.30\textwidth}
\textbf{C3 (matched):}
\begin{verbatim}
In what country is
Normandy...?
{
  "answer": "location",
  "location": "",
  "relation": ""
}
\end{verbatim}
\end{minipage}

\vspace{1.5em}

\begin{minipage}[t]{0.30\textwidth}
\textbf{C4 (passage + matched):}
\begin{verbatim}
The Normans gave
their name...
In what country is
Normandy...?
{
  "answer": "location",
  "location": "",
  "relation": ""
}
\end{verbatim}
\end{minipage}
\hfill
\begin{minipage}[t]{0.30\textwidth}
\textbf{C5 (random):}
\begin{verbatim}
In what country is
Normandy...?
{
  "answer": "cause",
  "effect": "",
  "cause": "",
  "mechanism": ""
}
\end{verbatim}
\end{minipage}
\hfill
\begin{minipage}[t]{0.30\textwidth}
\textbf{C6 (passage + random):}
\begin{verbatim}
The Normans gave
their name...
In what country is
Normandy...?
{
  "answer": "cause",
  "effect": "",
  "cause": "",
  "mechanism": ""
}
\end{verbatim}
\end{minipage}

\caption{The six conditions for a single question. The question is \textit{``In what country is Normandy located?''} and the passage is its SQuAD source, which states that Normandy is ``a region in France'' (both abbreviated above). The matched schema (C3, C4) is the location structure for the question's class; the random schema (C5, C6) is drawn from a non-matching class, here the causal structure. Each condition supplies the question with some combination of the passage and a schema.}
\label{fig:six_conditions}
\end{figure}

Each model is evaluated on the six conditions using a fixed stratified sample of 1,000 questions drawn class-proportionally from the SQuAD development set. All conditions use identical decoding parameters, with temperature set to 0 for deterministic output, so that differences are attributable only to the scaffold configuration. Prediction length is recorded alongside accuracy as a measure of verbosity.

\section{The Fine-Tuning Study}

The fine-tuning study addresses RQ5 through RQ7. The six-model study varies model size across its six models; this study holds size constant and instead varies how the structure is learned and delivered, isolating the effect of training from the effect of scale. It runs on a single model, Llama 3.2 3B in its base form, with a separate low-rank adapter trained for each experiment (LoRA, rank 16) at half precision. Each adapter is trained for three epochs with the answer masked from the loss, so the model learns only to produce the answer from the scaffold. A fresh base model is loaded for each experiment so the adapters do not contaminate one another, and decoding is greedy, matching the deterministic setting of the six-model study.

The study comprises three trained experiments:

\begin{itemize}
  \item \textbf{FT1} trains on the question-only and matched-structure conditions together, teaching the model the extractive format alongside the matched schema.
  \item \textbf{FT2} trains on the matched and random structures together, so the model sees both a correct and an incorrect schema for each question.
  \item \textbf{FT3} trains on the question-only condition alone, with no structure, as a format-only control.
\end{itemize}

Two further conditions are evaluated on the untuned model as comparisons for the trained experiments. A prose rendering expresses the same empty structure as a sentence rather than a schema, isolating the effect of delivery form. A structured passage reorganizes the full source passage into a labeled form without pre-extracting the answer, testing whether structure can carry content as well as the raw passage does.

A further attempt filled the schema from the passage and supplied it directly, but this could not be made to work without leaking the answer into the input, so it is reported only as a failed probe and not carried as an experiment.

These map to the final three questions. RQ5, whether training rescues the schema, compares each trained model's matched condition against the untuned baseline. RQ6, whether the delivery form matters, compares the schema against a prose rendering of the same structure. RQ7, whether structure can carry content, compares a passage reorganized into structure against the raw passage.

The eval sample of 1,000 questions is excluded from all training data, so no question the model is evaluated on appears during training. As this study uses a different model, runtime, and precision from the six-model study, the two are not compared on absolute values; only differences measured against each study's own baseline are compared across them.

\section{Models Evaluated}

The six-model study evaluates six open-source decoder models spanning a range of sizes and architectures, from 1.3B to 8.2B parameters: Llama 3.2 1B (1.3B), Gemma 2 2B (2.6B), Llama 3.2 3B (3.2B), Phi-3 Mini (3.8B), Llama 3 8B (8.0B), and Qwen 3 8B (8.2B). The two 8B models hold size fixed while varying architecture, additionally isolating architecture from scale. All six are served locally through Ollama in their default quantized form. Qwen 3 8B is a reasoning model that generates an intermediate reasoning trace by default. This was disabled so that it answers directly, on the same footing as the other models.

The fine-tuning study uses one of these models, Llama 3.2 3B, in its base (non-instruction-tuned) form. Training from the base model isolates what the adapter contributes: with no instruction tuning already in place, any gain is attributable to the training rather than to prior alignment.

\section{Evaluation Metrics}
\label{sec:metrics}

Model performance is measured by Exact Match (EM) and token-level F1, both standard metrics for SQuAD and directly comparable to published baselines. In the six-model study, results are reported as deltas between conditions: the matched structure against baseline (RQ1), the matched against the random structure (RQ2), and structure against the passage (RQ3), with all deltas compared across models (RQ4). Each reported per-condition mean carries its standard error of the mean (SEM), computed as the standard deviation of the per-question scores divided by $\sqrt{n}$ with $n = 1{,}000$, shown in parentheses in the corresponding tables.

The model produces free-form text rather than extractive spans. A raw EM comparison between the full output and the gold answer would therefore undercount correct responses: for example, ``Normandy is a region in France'' contains the correct answer ``France'' but fails string-level EM. A max-F1 substring extraction step is therefore applied. For each prediction, every contiguous token subsequence is evaluated against each gold answer and the span that maximizes token-level F1 is selected. EM is then computed on this extracted span. The search window is bounded at twice the length of the longest gold answer for efficiency, since the optimal span is unlikely to exceed this length. Both raw EM and extracted EM (EM\textsubscript{ext}) are reported; EM\textsubscript{ext} serves as the primary exact-match metric. This approach is prompt-agnostic: it does not alter experimental conditions or introduce confounding instruction tokens.

Token-level F1 is likewise reported in two forms: raw F1 (computed over the full model output) and extracted F1 (F1\textsubscript{ext}, computed on the extracted span). Raw F1 penalizes verbose outputs even when the correct answer is present, since precision drops with every extraneous token. F1\textsubscript{ext} isolates answer quality from verbosity and is used as the primary F1 metric alongside EM\textsubscript{ext}. Mean prediction length (token count) is reported separately as an explicit verbosity measure.

In the fine-tuning study, each reported difference carries a 95\% confidence interval computed by paired bootstrap resampling over the per-question F1\textsubscript{ext} scores. A difference is treated as significant only when its interval excludes zero, distinguishing differences that are reliably nonzero from those that are not.

% ============================================================
\chapter{Results}
\label{ch:results}
% ============================================================

\section{Overview}

A matched answer structure does not improve the accuracy of a small model on extractive question answering under any of the conditions tested. The taxonomy fails whether the structure is matched or random, whether or not the passage is present, across model sizes, and after fine-tuning. The single factor that changed accuracy was the form in which the structure was delivered, not its content. A prose rendering outperformed the imposed schema, and that difference persisted both when the structure carried the full passage content and after fine-tuning.

\section{The Negative Result Across Angles}

On extractive QA a matched answer structure does not improve accuracy for any model or in any condition tested. Table~\ref{tab:f1ext} reports F1\textsubscript{ext} for the six instruction-tuned models across conditions 1--6, and Table~\ref{tab:emext} reports exact match (EM\textsubscript{ext}) for the same conditions.

\begin{table}[h]
\centering
\small
\setlength{\tabcolsep}{3pt}
\begin{tabular}{lcccccc}
\toprule
Model & C1 base & C2 pass & C3 match & C4 p+m & C5 rand & C6 p+r \\
\midrule
Llama 3.2 1B & 0.335 (0.012) & 0.791 (0.011) & 0.299 (0.011) & 0.738 (0.012) & 0.300 (0.011) & 0.730 (0.012) \\
Gemma 2 2B   & 0.388 (0.012) & 0.919 (0.007) & 0.337 (0.012) & 0.865 (0.009) & 0.327 (0.012) & 0.804 (0.011) \\
Llama 3.2 3B & 0.369 (0.012) & 0.906 (0.008) & 0.343 (0.012) & 0.891 (0.008) & 0.360 (0.012) & 0.874 (0.009) \\
Phi-3 Mini   & 0.415 (0.012) & 0.834 (0.009) & 0.388 (0.012) & 0.808 (0.010) & 0.372 (0.012) & 0.767 (0.011) \\
Llama 3 8B   & 0.441 (0.012) & 0.930 (0.006) & 0.403 (0.012) & 0.915 (0.007) & 0.397 (0.012) & 0.892 (0.009) \\
Qwen 3 8B    & 0.482 (0.012) & 0.952 (0.005) & 0.384 (0.012) & 0.902 (0.008) & 0.357 (0.012) & 0.882 (0.009) \\
\bottomrule
\end{tabular}
\caption{F1\textsubscript{ext} by model and condition (1,000 questions). Token-level F1 computed on the max-F1 extracted span. Standard error of the mean (SEM) is shown in parentheses, computed over the per-question scores ($n = 1{,}000$). C1--C6 correspond to the conditions in \sref{sec:sixmodel-method}.}
\label{tab:f1ext}
\end{table}

\begin{table}[h]
\centering
\small
\setlength{\tabcolsep}{3pt}
\begin{tabular}{lcccccc}
\toprule
Model & C1 base & C2 pass & C3 match & C4 p+m & C5 rand & C6 p+r \\
\midrule
Llama 3.2 1B & 0.122 (0.010) & 0.646 (0.015) & 0.109 (0.010) & 0.594 (0.016) & 0.111 (0.010) & 0.589 (0.016) \\
Gemma 2 2B   & 0.166 (0.012) & 0.818 (0.012) & 0.121 (0.010) & 0.758 (0.014) & 0.123 (0.010) & 0.700 (0.014) \\
Llama 3.2 3B & 0.160 (0.012) & 0.801 (0.013) & 0.142 (0.011) & 0.785 (0.013) & 0.151 (0.011) & 0.766 (0.013) \\
Phi-3 Mini   & 0.195 (0.013) & 0.668 (0.015) & 0.160 (0.012) & 0.657 (0.015) & 0.160 (0.012) & 0.614 (0.015) \\
Llama 3 8B   & 0.215 (0.013) & 0.839 (0.012) & 0.183 (0.012) & 0.816 (0.012) & 0.187 (0.012) & 0.800 (0.013) \\
Qwen 3 8B    & 0.240 (0.014) & 0.868 (0.011) & 0.178 (0.012) & 0.823 (0.012) & 0.168 (0.012) & 0.820 (0.013) \\
\bottomrule
\end{tabular}
\caption{EM\textsubscript{ext} by model and condition (1,000 questions). Exact match computed on the extracted span. Standard error of the mean (SEM) is shown in parentheses, computed over the per-question scores ($n = 1{,}000$). C1--C6 correspond to the conditions in \sref{sec:sixmodel-method}.}
\label{tab:emext}
\end{table}

The two metrics agree on every comparison that follows: no scaffold condition improves over its no-structure counterpart under either metric. The negative result is therefore robust to the choice of metric and is not an artifact of the partial credit that F1 assigns. The discussion refers to F1\textsubscript{ext} for concreteness; the corresponding EM\textsubscript{ext} values follow the same ordering in Table~\ref{tab:emext}.

\paragraph{Matched structure and accuracy (RQ1).}
A matched structure does not improve accuracy in any model. On extractive QA, where answers are short spans with little internal structure for a schema to organize, the effect is consistently negative. Without a passage, the matched structure (C3) is lower than the question-only baseline (C1) for all six models. With a passage, the combined condition (C4) is lower than the passage-only condition (C2) for all six. The structure effect without a passage (C3 $-$ C1) is negative in every case, ranging from $-0.026$ on Llama 3.2 3B to $-0.098$ on Qwen 3 8B. Exact match follows the same pattern, with the drop from C2 to C4 reaching 0.060 on Gemma 2 2B.

\paragraph{Specificity of the structure (RQ2).}
The six-model study shows only a weak ordering in favor of the matched structure, and the fine-tuning study does not confirm it. With a passage, the matched condition (C4) exceeds the random condition (C6) for all six models, most clearly on Gemma 2 2B, at 0.865 against 0.804, and Phi-3 Mini, at 0.808 against 0.767. This ordering is fragile. In the fine-tuning study the matched schema was marginally worse than the random one without a passage, the opposite direction. On extractive QA the correctness of the structure therefore confers no reliable advantage, and the weak six-model ordering should not be read as evidence that the taxonomy supplies a signal the model can use.

\paragraph{Independence from the passage (RQ3).}
The structure does not contribute independently of the passage. The passage-only condition (C2) is the strongest or joint-strongest condition for every model, with F1\textsubscript{ext} between 0.791 and 0.952 when the passage is present against 0.335 to 0.482 when it is absent. Adding structure on top of the passage (C4) never exceeds the passage alone. On an extractive task the answer is contained in the passage and the task is to locate it, so the passage dominates, and the structure neither adds to it nor substitutes for it.

\paragraph{Consistency across model sizes (RQ4).}
No conclusive evidence of a size-dependent effect was found. In the six-model study the structure penalty (C3 $-$ C1) is negative for every model, but its magnitude does not track parameter count, and the two 8B models differ more from each other than either does from the 1B model. A size-dependent effect appears only in the additional flat and sentence format tests (\S\ref{sec:exploratory}), where the scaffold penalty on the 1B model is markedly larger than on the 8B model. If anything, this points toward smaller models being harmed more than larger ones, though the six-model study does not corroborate it. The question is revisited in \cref{ch:discussion}.

\paragraph{Verbosity.}
The scaffold's effect on output length depends on whether a passage is present, and in neither case does the change correspond to improved accuracy. Without a passage, adding structure shortens outputs; on Llama 3 8B the mean prediction length falls from 70.6 tokens at baseline to 57.6 with matched structure. With a passage, the effect reverses and structure lengthens outputs, from 27.4 tokens to 38.0 on the same model. The scaffold perturbs the form of the answer in both directions without improving its content.

\section{The Fine-Tuning Study}

Table~\ref{tab:finetune} reports F1\textsubscript{ext} for the untuned baseline and the three trained experiments (FT1--FT3, defined in \cref{ch:method}) across the conditions relevant to RQ5, RQ6, and RQ7. The untuned baseline is Llama 3.2 3B with no adapter, and every difference below is measured against it. Each difference is reported with a 95\% confidence interval, and is called significant only where that interval excludes zero.

\begin{table}[h]
\centering
\begin{tabular}{lcccc}
\toprule
Condition & Untuned & FT1 & FT2 & FT3 \\
\midrule
Question only        & 0.185 & 0.260 & 0.265 & 0.264 \\
Matched schema       & 0.128 & 0.255 & 0.258 & 0.228 \\
Random schema        & 0.149 & 0.254 & 0.266 & 0.218 \\
Prose structure      & 0.253 & 0.253 & 0.252 & 0.254 \\
\midrule
Passage              & 0.380 & --    & --    & --    \\
Structured passage   & 0.301 & --    & --    & --    \\
\bottomrule
\end{tabular}
\caption{Fine-tuning study on Llama 3.2 3B (F1\textsubscript{ext}, 1,000 questions). The untuned column is the same model with no adapter and anchors every paired difference. FT1--FT3 are the trained experiments defined in \cref{ch:method}: FT1 mixes no-structure and matched-schema pairs, FT2 mixes matched and random schemas, and FT3 trains on no-structure pairs only. The two conditions below the rule (passage and structured passage) are inference-time probes on the untuned model only, since the trained experiments used no passage and the passage is off-distribution for them. These values are not comparable to the six-model tables, which use a different model family, runtime, and precision.}
\label{tab:finetune}
\end{table}

\paragraph{Fine-tuning and the schema (RQ5).}
Training raises overall accuracy substantially, and the trained matched condition exceeds the untuned matched condition by 0.127 [95\% CI 0.101, 0.154] for FT1, 0.130 [0.104, 0.158] for FT2, and 0.100 [0.076, 0.126] for FT3, each significant. However, this gain reflects the model learning the extractive-QA format rather than learning to exploit the schema, since the question-only baseline rose by a comparable amount and the schema still does not beat the no-structure condition after training. The gap between the matched schema and the question-only baseline is neutralized rather than made positive, at $-0.006$ [$-0.018$, 0.007] for FT1 and $-0.007$ [$-0.020$, 0.006] for FT2, neither significant, and it remains significantly negative at $-0.037$ [$-0.056$, $-0.016$] for FT3. Training teaches the format of the answer, not the use of the structure.

\paragraph{The form of delivery (RQ6).}
The form in which the structure is delivered does change accuracy. Rendered as prose rather than as a schema, the same empty structure improves accuracy on the untuned model by 0.125 [95\% CI 0.103, 0.146], the clearest positive result in the study. The difference closes after fine-tuning, with the prose-over-schema difference no longer significant for FT1 at $-0.002$ [$-0.018$, 0.014] or FT2 at $-0.006$ [$-0.023$, 0.009]. This indicates that the schema penalty is learnable rather than fixed. The finding is bounded to one prose rendering against one JSON schema on extractive QA for these small models, and should not be read as prose being preferable in general.

\paragraph{Structure carrying content (RQ7).}
Reorganizing the same content into structure reduces accuracy. The raw passage was rewritten into an organized, labeled form without access to the question or answer, so that all content is preserved and none is pre-extracted. This lowers untuned accuracy relative to the raw paragraph by 0.078 [95\% CI 0.051, 0.106], significant. Imposing structure on the same full content costs roughly 0.08, which confirms the delivery-form result at the level of content and not only for an empty schema. A further attempt filled the schema from the passage and supplied it directly, but it could not be made to work without leaking the answer into the input. With a filled-schema content retention of 0.932, the filled schema functions as a container for the answer, so its 0.960 score reflects that leakage rather than an effect of form. It is therefore reported only as a failed probe rather than a result.

\paragraph{Distance to the passage.}
Even after training, the structure does not reach prose carrying content. The best trained matched condition, 0.258 on FT2, remains 0.122 below the untuned passage condition at 0.380. Structure, whether empty or trained, does not close the distance to the passage.

\section{Additional Experiments}
\label{sec:exploratory}

These additional tests were run to rule out the possibility that the negative result was an artifact of how the structure was delivered rather than of the structure itself. Each varied a single aspect of delivery: the format, its placement in the prompt, the generation pipeline, model scale, or the training regime of the model. Table~\ref{tab:exploratory} pairs each scaffold variant with the passage-only condition from the same run. The runs differ in model and sample size, so values are comparable only within a row and not across them; each row establishes direction rather than a comparable magnitude.

\begin{table}[h]
\centering
\begin{tabular}{llcccc}
\toprule
Delivery variant & Model & Scaffold F1 & Passage-only F1 & Direction & $n$ \\
\midrule
Structure in system prompt      & llama3:8b   & 0.233 & 0.429 & worse         & 30* \\
Two-pass fill-then-generate      & llama3:8b   & 0.824 & 0.895 & worse         & 130* \\
Filled schema in system prompt   & llama3:8b   & 0.849 & 0.939 & worse         & 30* \\
Flat keywords + passage          & llama3:8b   & 0.918 & 0.937 & worse         & 50 \\
JSON scaffold + passage          & llama3:8b   & 0.870 & 0.937 & worse         & 50 \\
Sentence + passage               & llama3:8b   & 0.794 & 0.972 & worse         & 50 \\
\bottomrule
\end{tabular}
\caption{Additional delivery variants, each paired with its own passage-only reference (F1\textsubscript{ext}). Entries marked * are preliminary partial runs. The scaffold and reference values on each row come from the same model and sample, so the within-row comparison is valid; values are not comparable across rows, as even the passage-only reference varies with the sample drawn.}
\label{tab:exploratory}
\end{table}

A separate rescoring test, stripping JSON syntax from predictions before scoring, changed F1\textsubscript{ext} by less than 0.002 on the full 1,000-question set, ruling out a measurement artifact from the scaffold format.

Across every delivery variant, the scaffold underperforms the passage alone. Placement made the effect worse, not better: moving the structure into the system prompt drove F1\textsubscript{ext} down to 0.233 against a 0.429 reference, as the accompanying instruction over-constrained the output. A two-pass pipeline, in which the model first filled the schema from the passage and then generated an answer from the filled schema, also fell below the passage alone (0.824 versus 0.895); extracting the passage into a schema discards information. Flat keywords were the least harmful format but still did not reach the passage baseline, and the natural-language sentence scaffold was worse than both. Each is a variant of an imposed schema, and each sits below the passage; the cleaner contrast between a schema and a prose rendering of the same structure is drawn out in the fine-tuning study (RQ6).

Measured without a passage, the scaffold penalty (scaffold F1 minus baseline) is larger on the smaller model: $-0.115$ for flat keywords and $-0.172$ for the sentence scaffold on Llama 3.2 1B, against $-0.025$ and $-0.090$ for the same formats on Llama 3 8B. The smaller model is harmed more by the imposed structure, not less.

The runs on pre-populated prompts and base-model behavior deserve particular attention. Two alternative runs of Llama 3 8B were carried out, each varying one property: one with the instruction-tuned model in raw completion mode without its chat template, and one with the base model under the same raw prompt. Together with the templated results from the six-model study, these give three configurations, all run on the full 1,000 questions and reported in Table~\ref{tab:rawtuning}.

\begin{table}[h]
\centering
\begin{tabular}{lccc}
\toprule
Condition & Instruct + template & Instruct + raw & Base + raw \\
\midrule
C1 question-only    & 0.441 & 0.396 & 0.304 \\
C2 passage          & 0.930 & 0.943 & 0.776 \\
C3 matched          & 0.403 & 0.276 & 0.188 \\
C4 passage + matched & 0.915 & 0.811 & 0.686 \\
C5 random           & 0.397 & 0.277 & 0.217 \\
C6 passage + random & 0.892 & 0.801 & 0.734 \\
\bottomrule
\end{tabular}
\caption{Llama 3 8B across prompt format and training (F1\textsubscript{ext}, 1,000 questions). The instruction-tuned model is shown with its chat template and in raw completion mode; the base model is the untuned weights under the same raw prompt.}
\label{tab:rawtuning}
\end{table}

Removing the chat template changes little on the passage condition, from 0.930 to 0.943, so the template is not the source of the passage's advantage. Removing instruction tuning changes it substantially, from 0.943 to 0.776 under the same raw prompt. The scaffold reduces accuracy in all three configurations; the instruction-tuned model in raw mode falls from 0.943 to 0.811 once the matched schema is added, in the same direction as the templated runs. The two models also differ in their response to a correct schema. The instruction-tuned model scores the matched schema above the random one in raw mode as it does with the template, while the base model reverses this, scoring the random schema higher. These results are interpreted in \cref{ch:discussion}.

\paragraph{Explicit instruction to use the structure.}
The main design supplied the schema as context but never instructed the model to use it. A probe on Llama 3.2 3B over a 100-question subsample added the instruction ``Use the following answer structure to organize your response'' to the matched-schema conditions. Without a passage, the instruction did not change the outcome: the instructed schema scored 0.352 in F1\textsubscript{ext} against 0.371 for the schema alone, both below the 0.384 question-only baseline. With a passage, the instructed condition rose to 0.976 against 0.876 without the instruction. The instructed outputs, however, copied passage text verbatim into the schema's fields rather than composing an answer. These results are interpreted in \cref{ch:discussion}.

\paragraph{Frontier model, zero-shot.}
The six conditions were also run against a frontier model, Claude Opus 5.5, zero-shot, on a 100-question subsample of the evaluation set. Table~\ref{tab:frontier} reports F1\textsubscript{ext} and EM\textsubscript{ext}. Without a passage, the matched structure scored 0.627 against a question-only baseline of 0.644, and the matched and random structures scored 0.627 and 0.608. With a passage, every condition exceeded 0.96, and adding structure to the passage did not raise accuracy beyond the passage alone. The model answered in fluent prose rather than terse spans, which depresses raw F1 and makes F1\textsubscript{ext} and EM\textsubscript{ext} the appropriate measures, as elsewhere. This run uses a different model family and runtime, a smaller sample, and the model's default sampling, since Opus 5.5 does not accept a fixed temperature; its values are not comparable in magnitude to the earlier tables.

\begin{table}[h]
\centering
\begin{tabular}{lcccccc}
\toprule
Metric & C1 base & C2 pass & C3 match & C4 p+m & C5 rand & C6 p+r \\
\midrule
F1\textsubscript{ext} & 0.644 & 0.963 & 0.627 & 0.983 & 0.608 & 0.967 \\
EM\textsubscript{ext} & 0.470 & 0.910 & 0.470 & 0.940 & 0.450 & 0.940 \\
\bottomrule
\end{tabular}
\caption{Claude Opus 5.5, zero-shot, on a 100-question subsample (F1\textsubscript{ext} and EM\textsubscript{ext}). The same six conditions and prompts as the six-model study. Values are not comparable in magnitude to the earlier tables, which use different models, a 1,000-question sample, and a fixed temperature. C1--C6 correspond to the conditions in \sref{sec:sixmodel-method}.}
\label{tab:frontier}
\end{table}

\section{Comparison with Reference Performance}

No directly comparable evaluation of significantly larger models exists for this setting. SQuAD 1.1 predates current large language models and is widely treated as saturated, and the extraction-based scoring used here (\sref{sec:metrics}) is not identical to the official metric, so a like-for-like comparison against larger models cannot be drawn from published results.

The available reference points are the original benchmark figures. Human performance on SQuAD 1.1 is reported at 82.3 EM and 91.2 F1 on the test set, and BERT-large, a 340M-parameter encoder fine-tuned specifically for the task, reaches 84.1 EM and 90.9 F1 on the development set \cite{devlin2019bert}. With the passage, the smaller models here reach these figures on exact match, from 80.1 on Llama 3.2 3B to 86.8 on Qwen 3 8B, without any task-specific fine-tuning. Exact match is the more reliable basis for this comparison, since F1\textsubscript{ext} applies a substring-extraction step (\sref{sec:metrics}) that the official metric does not, and the F1 values should therefore be read as indicative. Even read cautiously, several of the smaller models sit within the range of the reference systems.

Broader evaluations of large models on SQuAD have been conducted, most notably HELM \cite{liang2022helm}, which benchmarks models up to 530B parameters. These are not directly comparable to the results here: they adapt models through prompting rather than fine-tuning, and HELM itself cautions that its findings may not transfer to fine-tuned settings.

% ============================================================
\chapter{Discussion}
\label{ch:discussion}
% ============================================================

\section{Principal Findings}

A matched answer structure did not improve extractive QA accuracy for any model tested, and in most cases reduced it. What failed, however, was not the idea of supporting a small model with structure. It was the form in which that structure was supplied. Given the source passage, every model improved sharply over its question-only baseline. Given the same information as an explicit schema, accuracy fell. The support a small model can use and the support the scaffold provided are not the same thing.

The difference between them is delivery. When the same answer structure was expressed as prose rather than as a schema, the penalty disappeared, which places the obstacle in the form of the input rather than in the structure it carries. A small model does not reject structure. It reads structure only in the form it was trained on, and a rigid schema is not that form.

\section{The Six-Model Study}
\label{sec:sixmodel-disc}

The six-model study establishes the negative across its full range. Without a passage, the matched structure fell below the question-only baseline for every model; with a passage, adding the structure lowered accuracy rather than raising it. The effect was consistent across six architectures spanning 1B to 8B parameters. It is therefore a property of prompt-supplied structure on this task rather than a quirk of one model.

The size of the passage effect confirms the negative is genuine. When a passage was present, accuracy rose by nearly half under the same evaluation that registered no benefit from the scaffold. The evaluation is sensitive enough to register a change this large, so a smaller improvement from the scaffold would have shown up had one existed. The scaffold's failure reflects an absence of effect rather than an inability to observe it. This reasoning carries a caveat particular to the task. The passage effect is large partly because the task is extractive, with the answer lying in the passage. The comparison should therefore be read as confirming the measure is sensitive, not as a baseline for what the scaffold could have achieved. A separate check rules out a different artifact: stripping the JSON syntax from predictions before scoring changed accuracy by less than 0.002, so the schema's punctuation is not what depressed its score.

Model size did not change this. The structure penalty was negative for every model, but its magnitude did not track parameter count. The two 8B models sat at opposite ends of the range. Were capacity the limiting factor, the penalty would ease as models grew, and it did not. This runs against the original motivation for the scaffold. Smaller models, being more limited, were expected to gain the most from external structure, yet if anything they were harmed slightly more. The study found no systematic size trend. The range tested reaches only to 8B, so the conclusion extends no further. Within that range, the scaffold's failure is not a matter of scale.

The scaffold's effect on output length tells the same story from another angle. Adding the structure changed how much the model produced, shortening answers without a passage and lengthening them with one, but neither change improved accuracy. The scaffold reshaped the form of the output without improving its content, which is the pattern that runs through every result.

The frontier-model probe is consistent with this. Run zero-shot on Claude Opus 5.5 (\sref{sec:exploratory}), the matched structure failed to beat the question-only baseline, and the passage dominated as before. That a frontier model shows the same behavior strengthens the central finding. The obstacle lies in how the structure is delivered rather than in the capacity of the model that receives it. This remains a single small probe and does not extend the thesis beyond the small models it concerns. It does close off the easiest objection. The schema did not fail only because the models tested were too small to use it.

\section{The Fine-Tuning Study}

The fine-tuning study tested whether the schema fails only because the model was never trained to read it. Training raised accuracy across the board, but the gain did not depend on the schema. The question-only condition rose as much as the matched-schema condition, so what training taught was the format of the task, not the use of the structure, consistent with the account of fine-tuning as narrowing toward a task distribution (\sref{sec:finetune-lit}). After training, the matched schema still did not outperform no schema at all. The penalty was removed without a benefit taking its place.

Two conditions show that the obstacle is form rather than content. Delivered as prose rather than as a schema, the same empty structure improved the untuned model's accuracy by 0.125 [95\% CI 0.103, 0.146]. The structure helped once it stopped looking like a schema. The reverse test points the same way: reorganizing the full source passage into a labeled structure, with no content removed, lowered accuracy by 0.078 [95\% CI 0.051, 0.106]. Prose carrying no added structure beat structure carrying the whole passage. In both tests the content was fixed and only the form changed, which locates the effect in the form.

This points to what the base model already holds. A pretrained model with no instruction tuning has learned the statistical regularities of ordinary text. It reads a passage fluently because a passage is ordinary text, and it stumbles on a schema because a schema is not. Fine-tuning can teach the model to tolerate the unfamiliar form, which is why the penalty disappears. Tolerance is not comprehension. The model learns to stop being hurt by the schema, not to draw anything from it. The bias toward prose is not a preference the model acquires but the shape of what it learned to begin with.

A further probe asked whether instructing the model to use the schema changes this (\sref{sec:exploratory}). It does not. Without a passage, an explicit instruction to use the structure left the schema below the question-only baseline. The uninstructed schema fell short in the same way. The schema did not fail because the model ignored it. With a passage the instructed condition did score higher, but the gain is the answer-container effect seen in the filled-schema probe (\cref{ch:results}). The instruction drove the model to copy passage text verbatim into the schema's fields. The answer is a span already present in the passage, so the extraction metric rewards that copy. The schema did not help the model answer; it held the content without organizing it.

The passage does more than supply the answer's content; it acts as in-context material the model conditions on. In-context learning operates over ordinary text (\sref{sec:prompt}). A well-formed passage is therefore not only a source of the answer but context the model can read directly. A schema offers the same information, but not in a form the mechanism can use.

\section{The Taxonomy}

The scaffold rested on the assumption, set out in \sref{sec:taxhyp}, that a finite set of answer-structure types exists and is distinguishable to the model. The results give this assumption only weak support, and less than the six-model study alone would suggest. In the six-model study a matched structure did score higher than a random one, which appeared to show the model responding to the correctness of the structure rather than to its mere presence. The fine-tuning study did not reproduce this. Without training, the matched schema scored marginally lower than the random one, by 0.022 [95\% CI 0.006, 0.037], and training left the two level. The apparent signal from the six-model study does not survive as a reliable effect.

This does not condemn the taxonomy, but it does bound what the study can claim for it. The taxonomy came from a single extractive dataset, and building it was straightforward: a small set of answer-structure types covered the questions. Whether it captures a distinction the model can use is a separate question from whether it can be built, and the results speak only to the first. On extractive QA the distinction does not carry, because the answers are short spans with little structure for a taxonomy to organize. A task with longer, more compositional answers might give the taxonomy more to describe, and there its structure could matter more. That possibility is untested here.

\section{Imposed vs Understood Structure}

The passage already supplies what the scaffold was meant to provide, and in a form the model reads directly. Consider the question of which country Normandy lies in. The passage states in prose that Normandy is a region in France. The matched scaffold instead restates the target as the schema \texttt{\{answer: location, location: "", relation: ""\}}, a template the model must map onto the same content before it can use it. The difference is not only that prose is the more familiar form. A base model is trained on natural language. What it encodes from that training is the relational and syntactic structure of prose itself, not the surface strings. A rigid key-value schema is a far shallower structure than the one the model already holds. The scaffold therefore asks the model to work from an impoverished description of an answer whose fuller structure the prose already carries, and it competes with the passage rather than reinforcing it.

This mechanism was not measured directly. The experiments recorded accuracy across conditions, not how the model processes prose and schema internally. The mechanism is therefore an interpretation of the results rather than a direct measurement. It rests on two findings the study does establish: the passage is the strongest input for every model, and an explicit schema lowers accuracy even when it is correct. The base-model runs support it from the other side.

If the passage helps because the model was trained to read its form, a model without that training should use it less. Removing the chat template from the instruction-tuned model barely changed its accuracy on the passage, from 0.930 to 0.943. Removing instruction tuning lowered it sharply, to 0.776 under the same prompt. The ability to use the passage rests on training, not on prompt format. The scaffold, meanwhile, failed in every configuration. The instruction-tuned model lost accuracy when the matched schema was added even with its template stripped away, so the loss cannot be blamed on the framing a template adds. The two models also failed differently: the instruction-tuned model scored the matched schema above the random one, and the base model reversed this. The sensitivity to a correct structure is therefore acquired through instruction tuning, not supplied by the schema. Structure can help a small model, but only once the model has been trained to read it.

\section{What This Teaches Us}
\label{sec:teaches}

The study set out to give a small model structure and found it already had some. What the model lacked was not structure but the ability to read it in any form other than the one it was trained on. That narrowness, not an absence of structure, is what the scaffold ran into.

Why the model is narrow in this way follows from how it represents language. The account developed in \sref{sec:llm} and \sref{sec:scale} treats the model as an engine over a learned distribution of language, one that shapes its output from whatever the input makes probable. The scaffold was an attempt to work through exactly this machinery: an answer structure was meant to narrow the prediction toward a better answer. The results show that narrowing works only when that input already belongs to the distribution the model reads. Prose narrows the prediction because it is a sample from that distribution. A rigid schema does not, and no amount of correctness in the schema changes the form it takes, likely because the schema's own tokens, its braces and keys, sit outside the text the model learned from. Every result in the study fits this refinement. The passage helps and the schema does not. Prose delivery of the same structure helps where the schema failed. Training buys tolerance of the unfamiliar form without ever making it useful.

This account is not proven here. The experiments measured what helps a small model answer, not how it represents its input, so the picture is the most economical explanation of the pattern rather than a demonstrated mechanism. Read that way, the study adds a piece of evidence to an existing picture of small models as readers of a distribution. Within that picture it locates a practical boundary: structure helps these models only when it arrives in the form they were trained to read.

\section{A Note in Hindsight}

The finding invites a sharper reading of why the schema failed. It is not that these models cannot handle JSON. They parse it and produce it readily, because the format is familiar from training, and stripping the schema's syntax before scoring barely changed the result (\sref{sec:sixmodel-disc}), so the punctuation itself was never the obstacle. The results suggest they are not equipped to treat a schema built for a particular task as a guide to answering it. As the mechanism operates, a language model can read a schema in only one way, as a fixed format for laying out information based on the incoming tokens. It cannot use a bespoke structure whose particular fields guide how to answer a particular question.

The result should also be read against the setting that produced it. Extractive question answering is a small and controlled task. The answer is a span already present in the passage, so the passage was always likely to be the strongest input, and an imposed structure had little room to add anything. SQuAD was chosen deliberately for this practicality: it supplied the passages the framework required, a body of questions from which to derive the taxonomy, and gold answers to measure against, without the cost of building a dataset from scratch. The same feature that made it a practical choice, the answer sitting in the passage, is what made the passage so hard to compete with. A task whose answers must be composed rather than located would give structure more to do, and would test the delivery question more sharply than this setting could.

It is tempting to extend the finding further, to the claim that these models interface well with text and poorly with anything else. The evidence here does not support that leap. What a model can and cannot read depends on what it was trained on, and that training data is rarely known in full. Without it, there is no way to say which forms a given model is equipped to interface with and which it is not. This is less a limitation of the present study than a standing difficulty in the field: conclusions about what these models can process are bounded by how little is known about what they were trained to process. This study is no exception; it worked with models whose training data it could not inspect, and its reading of what they can and cannot use is constrained accordingly.

This is raised to place the result, not to diminish it. Content is most likely to dominate form precisely in an extractive setting. A generative task, in which the model must compose rather than locate an answer, would test the role of structure more directly, and is left to future work.

% ============================================================
\chapter{Limitations}
\label{ch:limitations}
% ============================================================

Several limitations bound the scope of these conclusions. The evaluation uses SQuAD, whose answers are short factual spans. Without the passage the model must generate the answer from memory rather than locate it, ideally still resulting in a brief span. Such answers have little internal structure, so a scaffold meant to lay out the form of an answer has little to organize. Whether the result holds for tasks with composite answers, where form carries more weight, remains untested. The model sample is also narrow. The six-model study spans open-weight decoders between 1B and 8B, and the fine-tuning study rests on a single 3B model. Together these cover only a small part of the space of models, and no large-scale models were included. Whether the relationship between structure and capacity changes at that scale is therefore left open. The taxonomy was not derived from an existing framework but induced from a single dataset by an LLM annotator applying a fixed rule set. A taxonomy built differently might capture a different signal. The study establishes that such a taxonomy can exist, not what its ideal form is. A final limitation is an extraction-based F1 that recovers the best-matching span from the model output. Being more forgiving than the official SQuAD metric, it inflates the absolute numbers, so the figures here support comparison across conditions more than comparison with published results.

% ============================================================
\chapter{Future Work}
\label{ch:future}
% ============================================================

The clearest extension is to tasks whose answers are composite rather than extractive. A SQuAD answer is a short span already present in the passage. It has little internal structure of its own, so an imposed structure had little to organize (\cref{ch:limitations}). This is the setting least favorable to the method, and the negative result should be read against it. Tasks whose answers must be composed rather than located might change that balance. Multi-step reasoning, causal explanation, and long-form generation all produce answers with several parts standing in a definite relation, which is what an answer structure is meant to describe. Applying the method to these tasks would test whether a matched structure helps where form carries real weight, and would show whether the failure seen here is specific to extractive answers or general to prompt-supplied structure.

A second direction concerns the form in which structure is delivered. The results locate the failure in the schema rather than in the structure it carries. Rendered as prose, the same structure improved accuracy where the schema did not (\sref{sec:teaches}). The question this leaves open is how prose delivery should be composed. It is not yet clear how much detail a prose structure needs, which cues carry the guidance, or how much surrounding context helps before it begins to distract. Studying these factors would shift the focus from whether delivery form matters, which the results establish, to how a structured prose delivery is best written. One concrete form this question takes is specialized syntactic structure, where the schema's dimensions are carried in the grammatical relations of well-formed language rather than in keys and braces. Such a delivery would sit between the rigid schema that failed here and free prose, keeping the explicit dimensions while presenting them in the form the model reads most readily. Broadening the taxonomy beyond a single extractive dataset would run alongside this, showing whether one set of answer structures serves across tasks or each task calls for its own.

% ============================================================
\chapter{Conclusion}
\label{ch:conclusion}
% ============================================================

This work asked whether a small language model benefits from an answer structure matched to the question before generation. Underlying this was the expectation that an explicit description of the answer's form would narrow the model's prediction space, and that a model constrained in capacity would benefit most. A taxonomy of answer structures was derived from SQuAD and examined through two studies. A six-model study, spanning 1B to 8B parameters, varied the matched structure, a random structure, and the source passage across conditions. A fine-tuning study, on a single model, varied how the structure was trained and delivered.

The structure did not improve accuracy. It failed to help any model and in most cases reduced performance. The passage, by contrast, improved every one. What a small model lacks on this task is not a description of the answer's form but the structure already carried by a well-formed passage, one it has been pretrained to interpret.

The same structure expressed as prose improved accuracy where the schema did not, and training a model on the schema removed its penalty without turning it into a benefit. A small model does not reject structure. It reads structure only in the form it was trained on, and a rigid schema is not that form. This points to how these models represent language: as a learned distribution over ordinary text, into which a well-formed passage falls and a schema does not.

The taxonomy proved less decisive than its construction suggested. It could be derived and applied, and the six-model study hinted that a matched structure was read differently from a random one. That signal did not survive the fine-tuning study, where the matched and random schemas performed alike. Model size was not decisive either, as the effect did not track parameter count across the range tested.

Whether inference-time structure can help a small model generalize more broadly remains open, and the answer here is bounded to its setting. On extractive question answering, where the answer is a short span with little form of its own, an imposed structure did not help. The question is worth pursuing on tasks whose answers must be composed, where there is more structure for such an approach to organize. What the study establishes is narrower and firmer. Structure can help a small model, but only in a form it already reads. In practice, that means a well-formed passage rather than an imposed schema.


\backmatter

% ============================================================
% Bibliography
% ============================================================
% CSU template default is unsrt; the pre-migration draft used alpha.
\bibliographystyle{unsrt}
\bibliography{sample}

% ============================================================
\appendix
% ============================================================

\chapter{Code and Data Availability}
\label{app:code}

The experiment code and data supporting this work are available in a public repository at \url{https://github.com/EricTheRed12/answer-structure-prompting}. The repository contains the evaluation pipeline, the prompt scaffold builder, the answer-structure taxonomy and its classification scripts, the fine-tuning notebook, and the saved results for both studies, together with a README describing how to reproduce each study. The committed result files include the per-question scores from which the confidence intervals and standard errors reported in \cref{ch:results} are computed.

\end{document}
